\section{Discussion and extensions}

The precision comparisons distinguish three features of the observed experiment: the availability of primary outcomes, the number of covariate categories across which ascertainment must be accommodated, and the remaining ascertainment deficit. Their relative importance changes with the arrival regime. The resulting interpretation is a comparison of attainable precision over the stated model classes, with each regime retaining its own conditions and quantifiers.

\subsection{Informative outcomes and adjustment complexity}

The effective sample size \(N=nq\) measures outcome availability at the guaranteed occupied-cell arrival floor. Because individual cells may have arrival probabilities exceeding that floor, it is a conservative scale for outcome information. Under \cref{obj:ass:rare-arrival-slice}, the matched comparison in \cref{obj:thm:matched-minimax-frontier} separates the sampling contribution \(N^{-1}\) from the category contribution
\[
\left\{\frac{d}{N\log(e+N)}\right\}^{2},
\]
with their sum truncated at a constant. Thus increasing outcome availability improves both components, whereas increasing the alphabet size raises the cost of estimating the aggregate contrast across sparsely observed cells. The comparison incorporates the memberships observed for individuals whose primary outcomes have yet to arrive.

Along sequences satisfying the rare-arrival restriction and \(N\to\infty\), \cref{obj:prop:phase-boundaries} gives two distinct measures of manageable adjustment complexity. Consistency corresponds to \(d=o\{N\log(e+N)\}\), while squared-error risk of order \(N^{-1}\) corresponds to \(d=O\{\sqrt N\log(e+N)\}\). These thresholds distinguish category growth compatible with learning the treatment effect from category growth compatible with preserving the effective-sample-size precision benchmark. At each fixed coverage level, \cref{obj:thm:connected-interval-frontier} gives the corresponding interpretation for uniformly honest connected intervals: optimal maximal expected length has the order of the square root of the matched squared-error rate.

Near complete ascertainment, the deficit \(1-q\) provides a complementary measure of difficulty. The observed-outcome contrast attains the upper bound \(4/n+(1-q)^2\) throughout the unrestricted model in \cref{obj:thm:unrestricted-near-complete-arrival-envelope}. For \(d\ge1024n^2\), \cref{obj:thm:unrestricted-large-alphabet-deficit-lower} establishes the matching order \(n^{-1}+(1-q)^2\). At complete arrival, \cref{obj:thm:unrestricted-complete-arrival-frontier} gives squared-error order \(n^{-1}\) uniformly over every alphabet size. Together, these results explain how sufficiently complete ascertainment supports precision governed by sampling variation even under arbitrary category growth.

The quantifiers distinguish a uniform guarantee from a guarantee along a specified dimension sequence. In \cref{obj:prop:unrestricted-uniform-near-complete-elbow}, a single risk constant and an eventual sample-size threshold apply simultaneously to every \(d\ge1\). Such an \(n^{-1}\) guarantee is equivalent to \(1-q_n=O(n^{-1/2})\). A specified sequence can satisfy a different sufficient condition through its category growth: at any fixed \(q\in(0,1/2)\), \cref{obj:prop:unrestricted-fixed-q-phase-boundaries} characterizes \(n^{-1}\) risk by \(d_n=O\{\sqrt n\log(e+n)\}\). The comparison constants may depend on the fixed floor and the category-growth multiplier. The near-complete criterion therefore describes precision uniformly over arbitrary alphabets, while the fixed-floor criterion describes precision along sequences with controlled adjustment complexity.

\subsection{Surrogates and prospective trial interpretations}

The surrogate serves as an observed variable for conditioning outcome arrival. Under \cref{obj:ass:arrival-mar}, arrivals represent the outcome distribution within treatment--baseline--surrogate cells; \cref{obj:ass:occupied-cell-arrival} supplies positive ascertainment in every occupied cell. Its realized value may depend on assignment through \cref{obj:ass:surrogate-consistency}. The target remains the population average treatment effect on the primary outcome in \cref{obj:def:ate-functional}, with the cell aggregation specified in \cref{obj:def:cell-functional}.

The constant-surrogate specialization provides a useful benchmark for this interpretation. When both potential surrogates equal zero, the arrival condition reduces to conditioning on treatment and baseline category. For fixed alphabet size, sequences in the rare-arrival slice with diverging effective sample size retain minimax squared-error order \(N^{-1}\) within this specialization, as established by \cref{obj:prop:fixed-d-effective-sample-reduction}. That benchmark isolates the precision scale associated with outcome availability in a model where the surrogate supplies a single category.

Clinical-trial applications can be interpreted prospectively through these conditions. At a prespecified analysis horizon, a binary primary endpoint would supply the outcome, an available short-term measurement would supply the binary surrogate, and recorded baseline categories would determine the adjustment alphabet. Applying the comparisons would require independent sampling, balanced randomized assignment, consistency, and the stated conditional arrival and occupied-cell positivity conditions. The delayed binary endpoints considered by \citet[Abstract]{MarschnerBecker2001} provide motivation for this interpretation. The use of short-term outcomes and baseline covariates for interim decisions in \citet[Abstract]{VanLancker2020} supplies related motivation. Here, these examples connect the statistical observation scheme to prospective trial settings, while the precision conclusions retain the conditions of the cited internal results.

Further questions concerning scope and research priorities are collected in \cref{sec:discussion-limitations-and-open-questions}.
